% ============================================================================
% LATEX REVISIONS FOR REVIEWER RESPONSE - CORRECTED VERSION
% All issues from review fixed
% Generated: 2025-12-27
% ============================================================================

% ----------------------------------------------------------------------------
% ISSUE #1: Clarify Ensemble Aggregation (Section 3.2)
% ----------------------------------------------------------------------------

% ADD TO SECTION 3.2 after "Evaluate coverage on held-out validation and test sets"

\subsection{Ensemble for Variance Quantification}

We train 50 independent models with random seeds 42--91 to quantify uncertainty
due to model initialization. For each seed $s$:

\begin{enumerate}
    \item Train a separate LightGBM classifier $M_s$ with random seed $s$
    \item Split the validation set 50/50 into calibration ($\mathcal{D}_{\text{cal}}^s$)
          and evaluation ($\mathcal{D}_{\text{eval}}^s$) sets
    \item Calibrate an individual conformal predictor $\text{CP}_s$ on $\mathcal{D}_{\text{cal}}^s$
    \item Evaluate coverage on $\mathcal{D}_{\text{eval}}^s$ (validation coverage)
          and test set (test coverage)
\end{enumerate}

We report mean $\pm$ standard deviation across 50 trials. This is \textit{not}
ensemble prediction (averaging model outputs)---each seed represents an
independent experimental trial to measure model variance.

% ----------------------------------------------------------------------------
% ISSUE #2: Enhanced Table 1 with Median/IQR
% ----------------------------------------------------------------------------

% REPLACE existing Table 1 with this enhanced version
% REQUIRES: \usepackage{multirow} and \usepackage{booktabs} in preamble

\begin{table*}[t]
\centering
\caption{Coverage Degradation Under COVID-19 Distribution Shift (50 model seeds).
High-variance tasks ($^*$) show severely skewed distributions where mean $\pm$ std
is misleading; median and IQR provide robust statistics.}
\label{tab:main_results}
\small
\begin{tabular}{@{}lcccccc@{}}
\toprule
\multirow{2}{*}{Task} & \multirow{2}{*}{Classes} &
\multicolumn{2}{c}{Val Coverage (\%)} &
\multicolumn{2}{c}{Test Coverage (\%)} &
\multirow{2}{*}{Cat} \\
\cmidrule(lr){3-4} \cmidrule(lr){5-6}
& & Mean$\pm$SD & Med(IQR) & Mean$\pm$SD & Med(IQR) & \\
\midrule
s-shipcond   & 45  & 93.5$\pm$1.0  & 93.5(1.4)  & 21.8$\pm$27.0 & 13.5(39.7) & SEV   \\
s-group      & 459 & 83.6$\pm$6.7  & 86.8(6.8)  & 12.4$\pm$32.3 & \textbf{0.5(0.7)}   & SEV$^*$ \\
s-payterms   & 137 & 90.8$\pm$0.6  & 90.6(0.7)  & 13.7$\pm$27.0 & \textbf{0.1(0.3)}   & SEV   \\
i-plant      & 35  & 92.0$\pm$0.9  & 91.7(1.3)  & 81.4$\pm$8.4  & 81.8(12.7) & ROB   \\
i-shippoint  & 69  & 91.2$\pm$0.8  & 91.1(0.9)  & 72.7$\pm$36.1 & 89.9(7.0)  & SEV$^*$ \\
s-incoterms  & 13  & 95.5$\pm$0.6  & 95.6(0.9)  & 87.0$\pm$8.0  & 92.3(17.3) & ROB   \\
i-incoterms  & 13  & 95.0$\pm$0.4  & 95.0(0.6)  & 83.7$\pm$9.9  & 82.1(19.8) & ROB   \\
s-office     & 25  & 99.9$\pm$0.0  & 99.9(0.0)  & 99.9$\pm$0.0  & 99.9(0.0)  & ROB   \\
\bottomrule
\end{tabular}
\vspace{2mm}

\raggedright
\footnotesize
SEV = Severe ($>$70\% drop), ROB = Robust ($<$15\% drop).
$^*$High model variance (std $>$ 30\%) indicates knife-edge regime where small
initialization changes lead to qualitatively different learned representations.
\end{table*}

% ----------------------------------------------------------------------------
% ISSUE #2b: Add to Discussion Section
% ----------------------------------------------------------------------------

% ADD NEW SUBSECTION in Section 7 (Discussion)

\subsection{High Model Variance as Early Warning Signal}

Tasks marked with $^*$ in Table~\ref{tab:main_results} exhibit extreme model
variance (std $>$ 30\%) across random seeds. Examining median values reveals
these are \textit{not normally distributed}:

\begin{itemize}
    \item \textbf{s-group}: Mean test coverage 12.4\%, but median \textbf{0.5\%}.
    Most models fail catastrophically, but a few outliers inflate the mean.

    \item \textbf{i-shippoint}: Mean 72.7\%, but median \textbf{89.9\%}.
    Most models maintain coverage, but failures deflate the mean.
\end{itemize}

This suggests these tasks are in a \textit{knife-edge regime}: small differences
in random initialization lead to qualitatively different learned representations.
Such high variance \textit{before distribution shift} is an early warning that
conformal calibration will be unreliable, as the model's behavior is fundamentally unstable.

\textbf{Practical Implication:} Practitioners should monitor model variance
across seeds during development. Tasks with coefficient of variation $>$ 50\%
may require ensemble approaches or architectural changes before deploying
conformal prediction.

% ----------------------------------------------------------------------------
% ISSUE #3: Resolve Jaccard Contradiction (Section 4.4)
% ----------------------------------------------------------------------------

% REPLACE paragraph in Section 4.4 starting with "Surprising finding: Both tasks..."

\textbf{Surprising finding:} Both tasks exhibited near-complete feature
distribution shift with $\sim$0\% Jaccard similarity across their
\textit{top-5 most important features} (mean Jaccard = 0.02 for sales-shipcond,
0.03 for sales-office from SHAP analysis), yet their coverage drops differed
by 700$\times$. This rules out feature value overlap as the primary mechanism
and points to a more subtle phenomenon: feature importance dynamics under
distribution shift.

While sales-office does have one stable feature (SALESORGANIZATION, Jaccard = 0.61),
this feature accounts for only 20\% of total SHAP importance (11.46 out of 57),
allowing the model to redistribute importance across multiple features post-shift.
In contrast, sales-shipcond's dominant feature (SALESDOCUMENT, Jaccard = 0.02)
accounts for 45\% of importance (9.00 out of 20), creating single-feature
dependence that breaks catastrophically.

% UPDATE Table 2 caption
% NOTE: Adjust \ref{sec:shap} to match your actual section label

\caption{Feature Overlap for Primary Features. Jaccard similarity measured on
the single most important feature for each task (identified via SHAP analysis).
Overall model robustness depends on \textit{all} top features, not just the
primary one---see Section 4.4 for analysis of top-5 feature dynamics.}

% ----------------------------------------------------------------------------
% ISSUE #5: Expand Related Work Section
% ----------------------------------------------------------------------------

% ADD TO SECTION 2 (Related Work)

\subsection{Distribution Shift Benchmarks}

Recent benchmarks systematically study distribution shift across domains.
WILDS~\cite{koh2021wilds} provides standardized evaluation for in-the-wild shifts,
while Shifts~\cite{malinin2021shifts} focuses on safety-critical applications.
Gulrajani and Lopez-Paz~\cite{gulrajani2020search} show that many ``robust''
methods fail under temporal shift. We focus specifically on how
\textit{conformal prediction guarantees} degrade under temporal shift.

\subsection{Conformal Prediction Under Distribution Shift}

Tibshirani et al.~\cite{tibshirani2019conformal} study conformal prediction
under covariate shift with known propensity scores. Podkopaev and
Ramdas~\cite{podkopaev2021distribution} develop distribution-free prediction
sets under label shift. Barber et al.~\cite{barber2023conformal} provide a
comprehensive treatment of conformal prediction beyond exchangeability.
We focus on \textit{temporal shift with feature staleness}---a distinct failure
mode where feature distributions have zero overlap between train and test.

Gibbs and Cand\`es~\cite{gibbs2021adaptive} propose Adaptive Conformal Inference
(ACI) for non-stationary settings, with extensions by Zaffran et
al.~\cite{zaffran2022adaptive} for time series. We show in Section 5.1
that these adaptive methods \textit{fail} when feature overlap is $\sim$0\%---this
is not a calibration problem but a fundamental data shift that requires retraining.

\subsection{Interpretability and Model Debugging}

Our finding that SHAP importance concentration predicts conformal failure
connects to work on feature attribution~\cite{lundberg2017unified} and model
debugging~\cite{adebayo2018sanity}. While SHAP is typically used for
post-hoc explanation, we show it can serve as a \textit{pre-deployment diagnostic}
for conformal prediction robustness.

% ----------------------------------------------------------------------------
% ISSUE #6: Reproducibility Details Appendix
% ----------------------------------------------------------------------------

% ADD AS APPENDIX A (or supplementary material)

\section{Reproducibility Details}
\label{app:reproducibility}

\subsection{Dataset Temporal Splits}

We use the rel-salt supply chain dataset with the following temporal splits
aligned to COVID-19 milestones:

\begin{itemize}
    \item \textbf{Training}: Before February 1, 2020 (pre-COVID baseline)
    \item \textbf{Validation}: February 1--July 1, 2020 (COVID onset; WHO
    declared pandemic March 11, 2020)
    \item \textbf{Test}: After July 1, 2020 (COVID peak; Delta variant emergence)
\end{itemize}

\subsection{LightGBM Hyperparameters}

All models trained with:

\begin{verbatim}
{
  'objective': 'multiclass',
  'boosting_type': 'gbdt',
  'num_leaves': 31,
  'learning_rate': 0.05,
  'feature_fraction': 0.8,
  'bagging_fraction': 0.8,
  'bagging_freq': 5,
  'num_boost_round': 500,
  'early_stopping_rounds': 50,
  'seed': [42, 43, ..., 91]  # 50 seeds
}
\end{verbatim}

These are standard LightGBM defaults. No task-specific tuning was performed to
avoid overfitting to the validation set used for conformal calibration.

\subsection{Conformal Prediction Setup}

\begin{itemize}
    \item \textbf{Method}: Adaptive Prediction Sets (APS)~\cite{romano2020classification}
    \item \textbf{Target coverage}: $(1-\alpha) = 0.9$ (90\%)
    \item \textbf{Calibration}: Random 50\% split of validation set
    \item \textbf{Quantile level}: $q = \min\left(\frac{\lceil(n+1)(1-\alpha)\rceil}{n}, 1\right)$
    where $n$ is calibration set size
\end{itemize}

For each test example with predicted probabilities $\hat{\pi}$, we construct
prediction set $\mathcal{C}(\hat{\pi})$ containing top classes until cumulative
probability exceeds quantile $q$.

\subsection{SHAP Computation}

\begin{itemize}
    \item \textbf{Method}: TreeExplainer~\cite{lundberg2020local}
    \item \textbf{Subsample size}: 10,000 samples per split (to reduce computation)
    \item \textbf{Aggregation}: Mean absolute SHAP values across all classes for
    multiclass tasks: $\text{SHAP}_f = \frac{1}{C} \sum_{c=1}^C |\phi_f^{(c)}|$
    \item \textbf{Feature ranking}: Based on mean $|\text{SHAP}_f|$ across samples
\end{itemize}

\subsection{Computational Resources}

\begin{itemize}
    \item \textbf{Hardware}: Standard CPU (no GPU required for LightGBM)
    \item \textbf{Parallel execution}: 8 cores for 50-seed ensemble
    \item \textbf{Wall-clock time}: $\sim$3--4 hours for full 8-task, 50-seed suite
    \item \textbf{Peak memory}: $<$8GB RAM per task
\end{itemize}

\subsection{Code Availability}

Code will be made available upon publication. Implementation built on the
RelBench framework~\cite{fey2023relbench}.

% ----------------------------------------------------------------------------
% NEW REFERENCES TO ADD
% ----------------------------------------------------------------------------

% ADD TO references.bib

@inproceedings{koh2021wilds,
  title={{WILDS}: A benchmark of in-the-wild distribution shifts},
  author={Koh, Pang Wei and Sagawa, Shiori and Marklund, Henrik and Xie, Sang Michael and Zhang, Marvin and Balsubramani, Akshay and Hu, Weihua and Yasunaga, Michihiro and Phillips, Richard Lanas and Gao, Irena and others},
  booktitle={International Conference on Machine Learning (ICML)},
  pages={5637--5664},
  year={2021},
  organization={PMLR}
}

@inproceedings{gulrajani2020search,
  title={In search of lost domain generalization},
  author={Gulrajani, Ishaan and Lopez-Paz, David},
  booktitle={International Conference on Learning Representations (ICLR)},
  year={2021}
}

@inproceedings{podkopaev2021distribution,
  title={Distribution-free uncertainty quantification for classification under label shift},
  author={Podkopaev, Aleksandr and Ramdas, Aaditya},
  booktitle={Conference on Uncertainty in Artificial Intelligence (UAI)},
  pages={844--853},
  year={2021},
  organization={PMLR}
}

@article{barber2023conformal,
  title={Conformal prediction beyond exchangeability},
  author={Barber, Rina Foygel and Cand{\`e}s, Emmanuel J and Ramdas, Aaditya and Tibshirani, Ryan J},
  journal={Annals of Statistics},
  volume={51},
  number={2},
  pages={816--845},
  year={2023},
  publisher={Institute of Mathematical Statistics}
}

@inproceedings{zaffran2022adaptive,
  title={Adaptive conformal predictions for time series},
  author={Zaffran, Margaux and F{\'e}ron, Olivier and Goude, Yannig and Josse, Julie and Dieuleveut, Aymeric},
  booktitle={International Conference on Machine Learning (ICML)},
  pages={25834--25866},
  year={2022},
  organization={PMLR}
}

@inproceedings{adebayo2018sanity,
  title={Sanity checks for saliency maps},
  author={Adebayo, Julius and Gilmer, Justin and Muelly, Michael and Goodfellow, Ian and Hardt, Moritz and Kim, Been},
  booktitle={Advances in Neural Information Processing Systems (NeurIPS)},
  volume={31},
  year={2018}
}

@article{malinin2021shifts,
  title={Shifts: A dataset of real distributional shift across multiple large-scale tasks},
  author={Malinin, Andrey and Band, Neil and Chesnokov, German and Gal, Yarin and Gales, Mark JF and Noskov, Alexey and Ploskonosov, Andrey and Prokhorenkova, Liudmila and Provilkov, Ivan and Raina, Vatsal and others},
  journal={arXiv preprint arXiv:2107.07455},
  year={2021}
}

% CORRECTED: Changed from @inproceedings to @article, booktitle to journal
@article{lundberg2020local,
  title={From local explanations to global understanding with explainable {AI} for trees},
  author={Lundberg, Scott M and Erion, Gabriel and Chen, Hugh and DeGrave, Alex and Prutkin, Jordan M and Nair, Bala and Katz, Ronit and Himmelfarb, Jonathan and Bansal, Nisha and Lee, Su-In},
  journal={Nature Machine Intelligence},
  volume={2},
  number={1},
  pages={56--67},
  year={2020},
  publisher={Nature Publishing Group}
}

% ----------------------------------------------------------------------------
% END OF LATEX REVISIONS - CORRECTED VERSION
% ----------------------------------------------------------------------------
